\documentclass[12pt,letterpaper]{article}
\usepackage[margin=1in]{geometry}
\usepackage{amsmath,booktabs,tabularx}
\usepackage[hidelinks]{hyperref}

\setlength{\parindent}{0pt}
\setlength{\parskip}{6pt}
\setlength{\emergencystretch}{2em}
\raggedbottom

\newcommand{\heading}[1]{\par\addvspace{10pt}\noindent\textbf{#1}\par\nobreak}
\newcommand{\repo}{https://github.com/billzhengmaozhang/COP3530p2}
\newcommand{\teamname}{Oracle (Group 94)}
\newcommand{\teammembers}{Zhengmao Zhang (\texttt{billzhengmaozhang})\\
  Jacob Lindland (\texttt{jdljake})\\
  My Pham (\texttt{phammy237})}
\newcommand{\demolink}{Link to be added when the recording is ready}

\begin{document}

\begin{center}
COP3530 Project 2\\[8pt]
\textbf{MiniDB}\\[8pt]
November 2026
\end{center}

\heading{Administrative Information}

\begin{tabularx}{\linewidth}{@{}l X@{}}
\textbf{Team name} & \teamname\\
\textbf{Team members} & \parbox[t]{\linewidth}{\teammembers}\\
\textbf{GitHub repository} & \href{\repo}{billzhengmaozhang/COP3530p2}\\
\textbf{Video demo} & \demolink
\end{tabularx}

\heading{1. Extended and Refined Proposal}

\textbf{Problem and motivation.}
MiniDB is an in-memory database for stock trades. It answers exact-ID queries and
compares two ways of finding the same trade: a custom separate-chaining hash table
and a custom B+ tree. Searching an unsorted collection can require inspecting every
row. An index avoids a full record scan, but it takes work to build and maintain.
The project asks how much that work costs and whether the answer changes with
dataset size, insertion order, or whether a queried ID exists.

\textbf{Implemented features.}
The terminal menu loads and saves trade CSV files, inserts a validated trade, and
looks up a trade by ID using the index selected by the user. Successful insertion
updates both indexes. Duplicate IDs are rejected and a failed CSV load preserves
the active dataset. A successful load replaces the active dataset. The benchmark
menu generates synthetic data, compares the two indexes in isolation, prints
median timings, and saves per-trial results as CSV.

\textbf{Data.}
Each row has a positive 64-bit trade ID, a ticker symbol, a price in integer cents,
a positive share quantity, and a nonnegative Unix timestamp in seconds. Storing
prices in cents preserves their exact values. The generator assigns contiguous
unique IDs and draws the remaining fields from fixed distributions using a seed.
It can emit rows in ascending ID order or shuffle those same rows. The repository
includes a small synthetic sample CSV. The default benchmark generates 100,000,
250,000, and 500,000 trades in both orders.

\textbf{Tools and data structures.}
MiniDB uses C++17 and CMake. Catch2 v3.5.4 and CTest support the B+ tree
tests. An append-only \texttt{std::vector} owns the full records. Both indexes
implement \texttt{IIndex} and store only a trade ID and a zero-based row position.
The position stays valid when the vector reallocates. Lookups return an optional
position so that row zero and a missing ID are distinct. The C++ standard library
provides vectors, smart pointers, binary search, pseudorandom generation, file
streams, and \texttt{std::chrono::steady\_clock} for timing.

\newpage
\heading{1. Extended and Refined Proposal (continued)}

\textbf{Hash table.}
The hash table starts with 16 buckets. A 64-bit integer mixer spreads keys across
the buckets. A vector within each bucket holds colliding entries, which are
searched by exact ID. Before an insertion would push the load factor above
$0.75$, the implementation doubles the bucket count and redistributes the entries.

\textbf{B+ tree.}
The default maximum is 32 children per internal node. Binary search over
separators chooses a child, with equal keys routed to the right. A second binary
search locates the ID in the leaf. Leaves hold sorted keys and their row positions,
and a link points to the next leaf. An overflowing leaf splits and copies its
right half's first key into the parent. An overflowing internal node promotes a
separator. Splits can propagate to the root and create a new root, keeping all
leaves at the same depth.

\textbf{Benchmark procedure.}
For each size and order, the runner times building an empty index, inserting
$\max(1,\lfloor n/10\rfloor)$ more keys, finding existing keys, and finding missing
keys. The lookup workloads each use 100,000 queries by default. The two indexes
receive identical keys and alternate which runs first. Each configuration uses
one warmup and five measured trials by default. The runner verifies returned
positions and compares checksums, with structural checks outside the timed
regions where applicable. The benchmark runs independently of the active database.

\textbf{Distribution of responsibility.}
The proposal assigned Zhengmao Zhang the hash table, collision and resize
testing, and hash benchmark integration. Jacob Lindland was assigned the B+
tree, its insertion and split tests, and tree benchmark integration. My Pham
was assigned record storage, data generation, CSV I/O, the menu, and the
benchmark runner. The repository records Jacob's implementation of the B+
tree and its Catch2 tests, and Zhengmao's integration of the CLI, database,
and benchmark components.

\heading{2. Analysis}

\textbf{Changes from the proposal.}
The proposal described generating synthetic CSV datasets. The benchmark menu
generates its datasets directly in memory, which lets users run the suite
without preparing input files and keeps CSV parsing outside the measurements.
The menu focuses on exact-ID lookup so the comparison uses the same operations
on both indexes. The B+ tree retains linked leaves for a possible range-query
extension. CSV export follows the design's write-then-rename approach and uses
a unique staging directory so separate saves cannot share a temporary file.

\newpage
\heading{2. Analysis (continued)}

\textbf{Worst-case running time.}
Let $n$ be the number of indexed records, $m$ the number of rows imported, $b$ the
number of bytes in the input file, and $q$ the number of lookups in a batch.
Assume a fixed B+ tree node capacity and the default hash-table configuration.

\begin{tabularx}{\linewidth}{@{}l l X@{}}
\toprule
\textbf{Operation} & \textbf{Worst case} & \textbf{Reason}\\
\midrule
Hash lookup & $O(n)$ & One bucket can contain all keys.\\
Hash insertion & $O(n)$ & Scanning a chain or rehashing can take linear time.\\
B+ tree lookup & $O(\log n)$ & Binary search at each level of a balanced tree.\\
B+ tree insertion & $O(\log n)$ & A split may propagate up one root-to-leaf path.\\
Database insertion & $O(n)$ & Hashing or growing the record vector can take linear time.\\
CSV load & $O(b+m^2)$ & Read and parse the file, then build both indexes.\\
CSV save & $O(n)$ & Validate and write each record once.\\
Dataset generation & $O(n)$ & Generate each row and optionally shuffle the vector.\\
\bottomrule
\end{tabularx}

Database lookup inherits the selected index's bound because resolving the row
position and copying a bounded-size record takes constant time. A well-distributed
hash gives expected $O(1)$ lookup and expected amortized $O(1)$ insertion.
Building from empty takes $O(n^2)$ in the worst case for hashing and
$O(n\log n)$ for the B+ tree. An incremental insert batch adds $\Theta(n)$ keys
and has the same bounds. The timed lookup batches take $O(qn)$ and
$O(q\log n)$, respectively. Outside timing, the hash structure check can also
take $O(n^2)$ because it scans earlier entries in each bucket for duplicates.
The record store and both indexes together use $O(n)$ space under the default
configuration.

\textbf{Measured results.}
Each cell shows hash / B+ tree median nanoseconds per operation from the
project benchmark summary. Build time is divided by the number of records built.

\begin{center}
\setlength{\tabcolsep}{4pt}
\begin{tabular}{@{}rrcccc@{}}
\toprule
Rows & Order & Build & Insert & Hit & Miss\\
\midrule
100,000 & Sequential & 181.68 / 47.84 & 78.97 / 50.41 & 17.20 / 102.96 & 15.97 / 13.09\\
100,000 & Shuffled   & 166.65 / 177.78 & 64.40 / 145.94 & 14.46 / 92.47 & 12.46 / 11.64\\
250,000 & Sequential & 181.72 / 49.10 & 113.86 / 58.10 & 19.78 / 132.10 & 21.17 / 14.85\\
250,000 & Shuffled   & 163.74 / 203.98 & 89.10 / 166.52 & 19.34 / 117.31 & 19.64 / 12.43\\
500,000 & Sequential & 200.08 / 51.08 & 121.56 / 55.03 & 26.89 / 162.36 & 23.23 / 14.85\\
500,000 & Shuffled   & 195.14 / 227.95 & 125.66 / 183.37 & 29.47 / 148.55 & 25.55 / 15.77\\
\bottomrule
\end{tabular}
\end{center}

\newpage
\heading{2. Analysis (continued)}

\textbf{Interpretation.}
The hash table was faster on successful lookups in all six configurations. At
500,000 shuffled records it took 29.47 ns per hit versus 148.55 ns for the tree.
The B+ tree built faster and handled the additional inserts faster for all three
sequential datasets. The hash table led both workloads for all three shuffled
datasets. Increasing IDs send B+ tree inserts to the rightmost leaf and append
within it. Shuffling sends inserts across different leaves and can require
shifting entries within a leaf. These differences are consistent with the
measurements.

The B+ tree was faster on misses in this run. The benchmark generates missing
IDs above the largest stored ID and the additional insertion range. That makes
all tree misses follow the same rightmost path, which can keep the visited nodes
in cache. The result describes this out-of-range workload. Misses spread across
the key space would exercise a different access pattern.

\heading{3. Reflection}

\textbf{Project experience and challenges.}
The project combined implementing the indexes with designing a way to compare
them. The shared interface made it possible to use one benchmark for both
implementations. The B+ tree exposed a subtle correctness issue to test: when a
leaf splits, searching for the promoted separator must still find the record in
the leaf on the right. Small node capacities made it possible to force leaf,
internal, and root splits with manageable test sequences. Another challenge was
distinguishing an indexing failure from an input failure. Keeping the old
database until a CSV import has passed validation gave the application a clear
rule for handling malformed rows and duplicate IDs.

\textbf{What we would change.}
We would add shared index tests and a CSV-to-lookup smoke test early so every
component has the same acceptance criteria. We would also benchmark absent IDs
between stored keys and archive each run with its revision and build settings.

\textbf{What each member learned.}
Zhengmao Zhang, Jacob Lindland, and My Pham each learned the importance of
agreeing on the architecture early, assigning clear responsibilities, and
communicating as work progresses. MiniDB's shared index interface illustrates
why those habits matter. Both indexes had to return a row position for the same
trade ID so the database could retrieve the record and the benchmark could
compare the results. Clear ownership gives each member a component to validate,
and regular communication helps catch mismatched assumptions before integration.

\heading{References}

[1] Oracle (Group 94), \emph{MiniDB Database Indexing Proposal}.

[2] MiniDB, \emph{Component Design and Interface Specification}.
\href{\repo/blob/main/docs/MiniDB_Design.md}{\texttt{docs/MiniDB\_Design.md}}.

[3] MiniDB, \href{\repo/tree/main}{implementation and tests}.

[4] MiniDB, \href{\repo/blob/main/results/submission-benchmark.csv}{archived benchmark trials}.

\end{document}
